\documentclass{zhvt-classic}

% 使用本地源流明体（GenRyuMin2 TW）按文件路径加载，不依赖系统字体名
% 版框（内框／界栏／书口侧栏／鱼尾）改用与版心印章同一印泥棕色：
%   印章图片 yshs.png 实色 = RGB(135,68,52) = HTML 874434
\usepackage{xcolor}
\usepackage{tikz}
% sealbrown 双注册：xcolor 供 TikZ/pgf，L3 供 cls 版框
\definecolor{sealbrown}{HTML}{874434}
\ExplSyntaxOn
\color_set:nnn { sealbrown } { HTML } { 874434 }
\ExplSyntaxOff

\zhvtset{
  grid color       = sealbrown,
  main font        = GenRyuMin2TW-SB,
  main font path   = C:/Users/Administrator/texmf/fonts/opentype/public/GenRyuMin2TW/,
jiazhu font      = GenRyuMin2TW-M,
  jiazhu font path = C:/Users/Administrator/texmf/fonts/opentype/public/GenRyuMin2TW/,
  % 页码框容纳 4 字：正常页「九十七」+前置空 1 格=4 行恰好；
  % 100+ 页「一百零一」共 5 行超出，末字「一/二/三」被框裁掉（2026-10-04 实测），故放宽到 5
  shukou page number chars = 5,
  % 书口书名框高从默认 6 字减到 4 字，把省出的 2 字高拨给 chapter 框，
  % 使 6 字书名+卷名放进 9 字高的 chapter 框后不再压线（2026-10-05）
  shukou title chars = 4,
}

\title{史记}[吕不韦列传]
\author{汉·司马迁}
\maker{闲叟漫记·重製}


% ===== 索隐述赞样式（2026-11-08 四次定稿 · 愚大人亲手改定参数）=====
% ⚠ 竖排前提：自定义字体家族必须带 RawFeature={vertical:+vert}；
%   simhei.ttf 加 +vhal 会整体加载失败静默回退，只写 +vert；simfang 同理。
\newCJKfontfamily{\shuzanfont}[Path=C:/Windows/Fonts/,RawFeature={vertical:+vert}]{simhei.ttf}
\newCJKfontfamily{\shuzanbody}[Path=C:/Windows/Fonts/,RawFeature={vertical:+vert}]{simfang.ttf}
% 题头：152x30pt 圆角 6pt 棕底白字；baseline (0,2) 使整框落入界行格子不压版框线；
% 白字节点 x 不补偿、y +8.2pt（像素级墨迹实测 dx/dy≈0，L/R 留白约 20pt）。
\newcommand{\shuzanhead}[1]{%
  \begin{tikzpicture}[baseline={(0,2)}]
    \fill[sealbrown, rounded corners=6pt]
      (-76pt,-15pt) rectangle (76pt,15pt);
    \node[text=white, inner sep=0pt, anchor=center] at (0pt,8.2pt)
      {\shuzanfont\fontsize{23}{23}\selectfont #1};
  \end{tikzpicture}%
}
% ★ 修正书口卷名位置：chapter 框起始位置改回 cls 默认值，避免全局影响。
%   只在每个卷首页通过 \juan 命令里的空命令微调该页位置（安全，不会跨页污染）。
\ExplSyntaxOn
\AtBeginDocument{
  \dim_set:Nn \l__zhvt_shukou_chapter_title_start_dim
    { \l__zhvt_shukou_title_height_dim + \l__zhvt_font_size_dim * 2 }
  \dim_set:Nn \l__zhvt_shukou_page_number_start_dim
    { \l__zhvt_shukou_chapter_title_start_dim + \l__zhvt_shukou_chapter_title_height_dim }
}
\ExplSyntaxOff

\begin{document}

% ★★ 版心（书口）排布覆写：把书名放到鱼尾之下（2026-10-04）
%
%   cls 默认版心自上而下：书名（title 静态框，在鱼尾之上）→ 鱼尾 → 篇名（chapter 动态框）→ 页码
%   本模板改为：          鱼尾（位置不动）→ 书名 → 篇名 → 页码分割线 → 页码
%
%   手法：把 cls 放在 title 静态框里的书名，改放进 chapter 动态框（鱼尾之下）；
%         title 静态框留空。书名取 \l__zhvt_title_str，即 \title{...} 的第一参数，
%         所以换书时只改 \title，不必在此重复写书名。
%
%   ⚠ 三条必须遵守的写法（2026-10-04 实测）：
%     ① 汉字必须用 \hbox{...} 包住。flowfram 会把框内容装进一个宽 14.4mm（= 一个 line_sep
%        = 书口栏宽）的纵向 minipage 并按该宽度折行，而整页靠 PDF /Rotate 90 显示、
%        框内文字本身不变向 —— 不包 \hbox 时「围炉夜话」4 字会被逐字折成 4 行，
%        在版面上就成了横排 4 字并横向溢出版框。
%     ② 竖排下**一个 \hbox 就是一整列**。书名与篇名要同处一列，必须放进同一个 \hbox；
%        分成两个 \hbox 时，后一个会被当成新的一列、甩到隔壁栏去
%        （实测「序」跑到书口栏外侧 14.4mm 的正文列上）。中间的 \quad 在框内是横向空白，
%        显示出来正好是书名与篇名之间的纵向间隔。
%     ③ ★ 内容盒要写成 **与框等宽的定宽盒** `\hbox to \l__zhvt_shukou_chapter_title_height_dim{…\hfil}`。
%        实测：chapter 框是把内容盒沿纵向**居中**摆放的，内容宽度一变，书名就跟着漂：
%        自然宽度的 \hbox（书名 38.4mm）落在 x=157.2mm，加篇名后（57.6mm）落在 x=152.4mm，
%        目录页（无篇名）与正文页的书名就差半格。做成等宽盒后，无论有没有篇名，
%        内容固定从框顶起排，实测两页均落在 x=142.8mm，不再漂。
\ExplSyntaxOn
% 把书口 chapter 框（书名+篇名）起始位置下移 1 字高，避免 4 字卷名顶到鱼尾
% ⚠ cls 的 \__zhvt_calculate_dimensions: 在 \AtBeginDocument 里执行，会重算该 dim；
%   必须把覆写挂到 begin document 之后，否则会被二次覆盖（2026-10-05 实测）。
% (旧方案已移到导言区，直接修改 chapter 框起始位置)
% 卷标题命令 #1=书口篇名(鱼尾) #2=正文卷名（含卷序，卷名内空格用 ~）
%
% ★★ 目录页两项修订（2026-10-05，愚大人指出）★★
%   ① 页码起点统一对齐：
%      原状 —— 目录项写成 \par\hspace{4em}卷名\hspace{4em}页码，页码起点 = 卷名实宽 + 4em。
%      卷名实宽随字数变：「卷一~梁惠王上」6 字 +~、「卷七~离娄上」5 字 +~，
%      差 1 个字位 → 卷七~卷十 4 条的页码整体比其余 10 条高 27.21pt（一个字），看着上下不齐。
%      改法 —— 卷名补全角空格（宽 = 1 字 = 27.215pt）补足到 **7 字符位**再写进目录，
%      14 条页码起点遂全部落在同一条竖线上（起点对齐；末字随页码字数自然参差）。
%      ⚠ 只对目录里的副本 \l_tmpc_tl 补，正文卷名仍用 \l_tmpa_tl 原样 —— 正文像素不受影响。
%   ② 卷名整体上移半字位：前导 \hspace 由 4em 改 3.5em。
%      实测原状目录块在版框内偏下（天头留白 139.6pt / 地脚留白 110.7pt），
%      上移 0.5em = 13.6pt 后上下基本均分，视觉居中。
\NewDocumentCommand { \juan } { mm } {
	\clearpage
	\tl_set:Nn \l__zhvt_chapter_title_tl { #1 }
	\tl_set:Nn \l_tmpa_tl { #2 }
	\tl_set_eq:NN \l_tmpc_tl \l_tmpa_tl
	\int_step_inline:nn { 7 - \str_count:n { #2 } }
	  { \tl_put_right:Nn \l_tmpc_tl { 　 } }
	\tl_set:Nx \l_tmpb_tl { \zhnum { page } }
	\tl_put_left:Nn \l_tmpb_tl { \hspace { 4em } }
	\tl_put_left:Nn \l_tmpb_tl { \l_tmpc_tl }
	\tl_put_left:Nn \l_tmpb_tl { \par \hspace { 3.5em } }
	\addtocontents { toc } { \l_tmpb_tl }
	\skip_horizontal:n { \l__zhvt_font_size_dim * 4 }
	\l_tmpa_tl
	\par
	\thispagestyle { empty }
}
\cs_set_protected:Npn \__zhvt_set_shukou_content: {
	% ★ 鱼尾上方留白（title 静态框）嵌入闲叟钤印：
	%   整页 /Rotate 90，故此处把图用 \rotatebox{90} 转正为竖章；
	%   尺寸卡在版心条幅内（书内约 7mm 宽 × 33.8mm 高，留余量不压鱼尾）。
	\setstaticcontents * { title } {
		\vbox_to_ht:nn { \l__zhvt_baseline_skip_dim } {
			\vskip 14.2pt
			\hbox_to_wd:nn { \l__zhvt_shukou_title_height_dim } {
				\hfil
				\rotatebox [ origin = c ] { 90 } { \includegraphics [ width = 7mm , height = 33.8mm ] { yshs.png } }
				\hfil
			}
		}
	}
% ★ 版心（书口）统一改法（2026-11-16 愚大人定）：18pt＋框内居中＋书口文字同版框色
%   ① 原 27.2pt 时 7 字长篇名（如「廉颇蔺相如列传」）末字溢出 chapter 框、压住页码；
%      改 \fontsize{18pt}{18pt}，前后 \hfil 令字串在框内居中（不随篇名字数漂移）。
%   ② \vbox{\vskip 8.2pt} 定量下移，抵消缩字号后栏宽方向 8.2pt 上偏（像素实测 y 中心 543.3 = 原版）。
%      ⚠ \raisebox / \vspace* 在 flowfram 框内无效；换字号须同步换补偿：16pt→10.1、19pt→7.4。
%   ③ chapter 与 page~number 均置 \color{sealbrown}（HTML 874434），与版框／界行同色；正文仍黑。
%   ⚠ 页码字位、钤印、鱼尾、界行、正文一律未动；五部合订本同此口径。
	\setdynamiccontents * { chapter }
	{ \color{sealbrown} \vbox { \vskip 8.2pt \hbox to \l__zhvt_shukou_chapter_title_height_dim
		{ \hfil \fontsize{18pt}{18pt}\selectfont \l__zhvt_title_str \quad \l__zhvt_chapter_title_tl \hfil } } }
	\setdynamiccontents * { page~number } { \color{sealbrown} \quad \zhnum { page } }
}
\__zhvt_set_shukou_content:
\ExplSyntaxOff

\maketitle

% 单篇，无多卷目录，不渲染目录页

\juan{吕不韦列传}{卷八十五~吕不韦列传第二十五}
\hspace*{2em}吕不韦者。阳翟大贾人也。往来贩贱卖贵。家累千金。\\
\hspace*{2em}秦昭王四十年。太子死。其四十二年。以其次子安国君为太子。安国君有子二十馀人。安国君有所甚爱姬。立以为正夫人。号曰华阳夫人。华阳夫人无子。安国君中男名子楚。子楚母曰夏姬。毋爱。子楚为秦质子於赵。秦数攻赵。赵不甚礼子楚。\\
\hspace*{2em}子楚。秦诸庶孽孙。质於诸侯。车乘进用不饶。居处困。不得意。吕不韦贾邯郸。见而怜之。曰﹁此奇货可居﹂。乃往见子楚。说曰。﹁吾能大子之门。﹂子楚笑曰。﹁且自大君之门。而乃大吾门。﹂吕不韦曰。﹁子不知也。吾门待子门而大。﹂子楚心知所谓。乃引与坐。深语。吕不韦曰。﹁秦王老矣。安国君得为太子。窃闻安国君爱幸华阳夫人。华阳夫人无子。能立適嗣者独华阳夫人耳。今子兄弟二十馀人。子又居中。不甚见幸。久质诸侯。即大王薨。安国君立为王。则子毋几得与长子及诸子旦暮在前者争为太子矣。﹂子楚曰。﹁然。为之柰何。﹂吕不韦曰。﹁子贫。客於此。非有以奉献於亲及结宾客也。不韦虽贫。请以千金为子西游。事安国君及华阳夫人。立子为適嗣。﹂子楚乃顿首曰。﹁必如君策。请得分秦国与君共之。﹂\\
\hspace*{2em}吕不韦乃以五百金与子楚。为进用。结宾客。而复以五百金买奇物玩好。自奉而西游秦。求见华阳夫人姊。而皆以其物献华阳夫人。因言子楚贤智。结诸侯宾客遍天下。常曰﹁楚也以夫人为天。日夜泣思太子及夫人﹂。夫人大喜。不韦因使其姊说夫人曰。﹁吾闻之。以色事人者。色衰而爱弛。今夫人事太子。甚爱而无子。不以此时蚤自结於诸子中贤孝者。举立以为適而子之。夫在则重尊。夫百岁之後。所子者为王。终不失势。此所谓一言而万世之利也。不以繁华时树本。即色衰爱弛後。虽欲开一语。尚可得乎。今子楚贤。而自知中男也。次不得为適。其母又不得幸。自附夫人。夫人诚以此时拔以为適。夫人则竟世有宠於秦矣。﹂华阳夫人以为然。承太子间。从容言子楚质於赵者绝贤。来往者皆称誉之。乃因涕泣曰。﹁妾幸得充後宫。不幸无子。原得子楚立以为適嗣。以讬妾身。﹂安国君许之。乃与夫人刻玉符。约以为適嗣。安国君及夫人因厚餽遗子楚。而请吕不韦傅之。子楚以此名誉益盛於诸侯。\\
\hspace*{2em}吕不韦取邯郸诸姬绝好善舞者与居。知有身。子楚从不韦饮。见而说之。因起为寿。请之。吕不韦怒。念业已破家为子楚。欲以钓奇。乃遂献其姬。姬自匿有身。至大期时。生子政。子楚遂立姬为夫人。\\
\hspace*{2em}秦昭王五十年。使王齮围邯郸。急。赵欲杀子楚。子楚与吕不韦谋。行金六百斤予守者吏。得脱。亡赴秦军。遂以得归。赵欲杀子楚妻子。子楚夫人赵豪家女也。得匿。以故母子竟得活。秦昭王五十六年。薨。太子安国君立为王。华阳夫人为王后。子楚为太子。赵亦奉子楚夫人及子政归秦。\\
\hspace*{2em}秦王立一年。薨。谥为孝文王。太子子楚代立。是为庄襄王。庄襄王所母华阳后为华阳太后。真母夏姬尊以为夏太后。庄襄王元年。以吕不韦为丞相。封为文信侯。食河南雒阳十万户。\\
\hspace*{2em}庄襄王即位三年。薨。太子政立为王。尊吕不韦为相国。号称﹁仲父﹂。秦王年少。太后时时窃私通吕不韦。不韦家僮万人。\\
\hspace*{2em}当是时。魏有信陵君。楚有春申君。赵有平原君。齐有孟尝君。皆下士喜宾客以相倾。吕不韦以秦之彊。羞不如。亦招致士。厚遇之。至食客三千人。是时诸侯多辩士。如荀卿之徒。著书布天下。吕不韦乃使其客人人著所闻。集论以为八览。六论。十二纪。二十馀万言。以为备天地万物古今之事。号曰吕氏春秋。布咸阳市门。悬千金其上。延诸侯游士宾客有能增损一字者予千金。\\
\hspace*{2em}始皇帝益壮。太后淫不止。吕不韦恐觉祸及己。乃私求大阴人嫪毐以为舍人。时纵倡乐。使毐以其阴关桐轮而行。令太后闻之。以啗太后。太后闻。果欲私得之。吕不韦乃进嫪毐。诈令人以腐罪告之。不韦又阴谓太后曰。﹁可事诈腐。则得给事中。﹂太后乃阴厚赐主腐者吏。诈论之。拔其须眉为宦者。遂得侍太后。太后私与通。绝爱之。有身。太后恐人知之。诈卜当避时。徙宫居雍。嫪毐常从。赏赐甚厚。事皆决於嫪毐。嫪毐家僮数千人。诸客求宦为嫪毐舍人千馀人。\\
\hspace*{2em}始皇七年。庄襄王母夏太后薨。孝文王后曰华阳太后。与孝文王会葬寿陵。夏太后子庄襄王葬芷阳。故夏太后独别葬杜东。曰﹁东望吾子。西望吾夫。後百年。旁当有万家邑﹂。\\
\hspace*{2em}始皇九年。有告嫪毐实非宦者。常与太后私乱。生子二人。皆匿之。与太后谋曰﹁王即薨。以子为後﹂。於是秦王下吏治。具得情实。事连相国吕不韦。九月。夷嫪毐三族。杀太后所生两子。而遂迁太后於雍。诸嫪毐舍人皆没其家而迁之蜀。王欲诛相国。为其奉先王功大。及宾客辩士为游说者众。王不忍致法。\\
\hspace*{2em}秦王十年十月。免相国吕不韦。及齐人茅焦说秦王。秦王乃迎太后於雍。归复咸阳。而出文信侯就国河南。\\
\hspace*{2em}岁馀。诸侯宾客使者相望於道。请文信侯。秦王恐其为变。乃赐文信侯书曰。﹁君何功於秦。秦封君河南。食十万户。君何亲於秦。号称仲父。其与家属徙处蜀。﹂吕不韦自度稍侵。恐诛。乃饮酖而死。秦王所加怒吕不韦。嫪毐皆已死。乃皆复归嫪毐舍人迁蜀者。\\
\hspace*{2em}始皇十九年。太后薨。谥为帝太后。与庄襄王会葬茝阳。\\
\hspace*{2em}太史公曰。不韦及嫪毐贵。封号文信侯。人之告嫪毐。毐闻之。秦王验左右。未发。上之雍郊。毐恐祸起。乃与党谋。矫太后玺发卒以反蕲年宫。发吏攻毐。毐败亡走。追斩之好畤。遂灭其宗。而吕不韦由此绌矣。孔子之所谓﹁闻﹂者。其吕子乎。\\

% ===== 索隐述赞（司马贞《史记索隐》）· 统一样式 =====
% 题头「索隐述赞曰」五字并入反白框；正文每四句一列，定宽盒 \hfil 居列中；
% \fontsize{24}{41} 列距对齐全书界行（121+41k）；\clearpage 使述赞独占一页。
\clearpage
\shuzanhead{索隐述赞曰}\linebreak
{\shuzanbody\fontsize{24}{41}\selectfont
\hbox to \hsize{\hfil 不韦钓奇。委质子楚。华阳立嗣。邯郸献女。\hfil}\linebreak
\hbox to \hsize{\hfil 及封河南。乃号仲父。徙蜀惩谤。悬金作语。\hfil}\linebreak
\hbox to \hsize{\hfil 筹策既成。富贵斯取。\hfil}%
\par}
\end{document}
